\documentclass[10pt,a4paper]{amsart}
\usepackage[margin=19mm]{geometry}
\usepackage{amsmath,amssymb,mathtools}
\usepackage[scr=rsfs,cal=euler]{mathalfa}
\usepackage{xfrac}
\usepackage[unicode,hidelinks,bookmarks=false]{hyperref}
\newcommand{\ie}{i.e.\ }
\newcommand{\cf}{cf.\ }
\newcommand{\resp}{respectively\ }
\newcommand{\Subsection}{\subsection}
\providecommand{\nobreakdash}{\nobreak\hbox{-}}
\setlength{\emergencystretch}{2em}
\multlinegap0pt
\usepackage{bm}
\usepackage{bbm}
\usepackage{dsfont}
\usepackage{xspace}
\usepackage{cancel}
\usepackage{mathtools}
\usepackage{amsthm}
\makeatletter
\@ifundefined{lemma}{\newtheorem{lemma}{Lemma}}{}
\makeatother
\def\Li{{\mathrm{Li}}}
\def\d{{\mathrm{d}}}
\title[Primes of the form $ax+by$ in certain intervals with small s]{Primes of the form $ax+by$ in certain intervals with small solutions}
\author{Yuchen Ding, Takao Komatsu, Honghu Liu}
\date{Source: arXiv:2510.01781v1 (CC BY 4.0)}
\begin{document}
\maketitle

\noindent\emph{Source and licence.} Primes of the form $ax+by$ in certain intervals with small solutions. Version arXiv:2510.01781v1 (CC BY 4.0), distributed under the Creative Commons Attribution 4.0 International licence, as recorded in the arXiv licence metadata for this exact version and shown on the abstract page https://arxiv.org/abs/2510.01781v1. Full rights notice: licenses/arxiv-2510.01781-CC-BY-4.0.txt.\par\smallskip

\noindent\textit{Editorial scope.} Counted unit: the Lemma whose source label is \texttt{lem-2} in the article (source0.tex lines 288--294, source label \texttt{lem-2}), delivered together with the article's own complete original proof of it (source0.tex lines 296--358, one \texttt{proof} environment of 63 lines). No mathematics is written, repaired or re-derived: statement and proof are the article's own text line for line. Required context retained uncounted: lem-1 (lines 281--285). Every definition, display and label of those lines is retained unchanged. Omitted from this package and not redistributed: 1--280, 286--287, 295--295, 359--1148. Nothing inside a retained excerpt is omitted or altered: every retained body is delivered exactly as the article's lines are written, so no omission receipt of any kind accompanies this package. Citations: the retained excerpts carry no citation command at all, so no bibliography entry of the article is transcribed and no reference is redistributed. Reference adaptation: none was needed and none was made; every reference inside the delivered unit resolves inside it, so no unresolved reference or citation is produced. Printed slips kept literal: no printed word, symbol, hypothesis or number of the counted statement or of its proof was altered. Layout and class: the article's own class is \texttt{amsart} with packages including amsmath, amsfonts, amssymb, amsthm, latexsym, enumerate, url, cases, graphicx, mathrsfs, hyperref, natbib, color, etoolbox; the wrapper is the lane's pinned \texttt{amsart} editorial wrapper at 10pt on A4, which restates the theorem-like environments the retained text needs and adds the article's own macro definitions that the retained text needs (\texttt{\textbackslash Li}, \texttt{\textbackslash d}), with extra wrapper packages (bm, bbm, dsfont, xspace, cancel, mathtools). The delivered numbering is the unit's own. Assets: no retained span contains a figure environment, a graphics-inclusion command, a PDF-page inclusion, a table or a TikZ picture; no image file is copied into this package, the package declares no asset dependency and no image file is redistributed with it. Licence: version arXiv:2510.01781v1 of this article is distributed under the Creative Commons Attribution 4.0 International licence, as recorded in the arXiv licence metadata for this exact version and shown on the abstract page https://arxiv.org/abs/2510.01781v1. A full rights notice is delivered at licenses/arxiv-2510.01781-CC-BY-4.0.txt. Characters: every retained excerpt is pure ASCII, so the delivered engine, XeLaTeX with its default Latin Modern text font, reports no missing character.
 \begin{lemma}\label{lem-1} For $x\ge 1474279333$, we have
 	$$\big|\pi (x)-\Li (x)\big|<0.0008375\frac{x}{\log^2 x}, $$
 	where
 	$$\Li (x)=\int_2^x \frac{\d t}{\log t} .$$
 \end{lemma}

 \begin{lemma}\label{lem-2}
 	Let $\ell$ be a positive integer and $x\geq \max\{\ell e^{2(\ell+1)},7\cdot 10^{5}\}$. Then
 	\begin{equation*}
 		\pi(x)-\pi(cx)\geq 0.11604925 \cdot \frac{1}{\ell+1} \frac{x}{\log x},
 	\end{equation*}
 	where $c=1-\frac{1}{8(\ell+1)}$.
 \end{lemma}

 \begin{proof}
 The proofs will be separated into two cases.
 	
 	\textbf{Case 1.}
 	For $x\geq \max\{\ell e^{2(\ell+1)},1.573\cdot 10^{9}\}$, we have 
 	$$
 	cx\geq \left( 1-\frac{1}{16} \right)x\geq \left( 1- \frac{1}{16}\right)\cdot 1.573\cdot 10^{9} > 1474279333.
 	$$
 	By Lemma \ref{lem-1}, we have
 	\begin{align}\label{lem2-ineq-1}
 		\pi(x)-\pi(cx)
 		&> \Li(x)-0.0008375\frac{x}{(\log x)^{2}}-\Li(cx)-0.0008375\frac{cx}{\big(\log (cx)\big)^{2}}\nonumber\\
 		&=\int_{cx}^{x}\frac{\d t}{\log t}-0.0008375\frac{x}{(\log x)^{2}}-0.0008375\frac{cx}{(\log c+\log x )^{2}}\nonumber\\
 		&\geq (1-c)\frac{x}{\log x}-0.0008375\frac{x}{(\log x)^{2}}-0.0008375\frac{cx}{(\log c+\log x)^{2}}\nonumber\\
 		&= \left( (1-c)-0.0008375\frac{1}{\log x}-0.0008375 \frac{c\log x}{(\log c+\log x)^{2}}\right)\frac{x}{\log x}\nonumber\\
 		&= \left( \frac{1}{8(\ell+1)}-\frac{0.0008375}{\log x}-\frac{0.0008375 c}{(\log c)^{2}/\log x+2\log c+\log x}\right)\frac{x}{\log x}.
 	\end{align}
 	Recalling that $c=1-\frac{1}{8(\ell+1)}$ and $x\geq \max\{\ell e^{2(\ell+1)},1.573\cdot 10^{9}\}$, we have 
 	\begin{align*}
 		\frac{(\log c)^{2}}{ \log x}+2\log c+ \log x
 		>2\log\left(1-\frac{1}{16}\right)+\log 10^{9}
 		>20
 		> 2 (\ell+1), \quad (1\leq \ell \leq 8)
 	\end{align*}
 	and
 	\begin{align*}
 		\frac{(\log c)^{2}}{\log x}+2\log c+ \log x 
 		&> 2\log \left(1-\frac{1}{8\cdot (9+1)}\right)+\log\left( \ell e^{2(\ell+1)} \right)\\
 		&>-0.025157565 +\log \ell+2(\ell+1)\\
 		&>2(\ell+1), \quad (\ell\geq 9).
 	\end{align*}
 	So, for any positive integer $\ell$ we have
 	\begin{equation}\label{lem2-ineq-2}
 		\frac{c}{(\log c)^{2}/\log x+2\log c+\log x}< \frac{1}{2} \frac{c}{ \ell+1}<\frac{1}{2}\frac{1}{\ell+1}.
 	\end{equation}
 	Hence, by (\ref{lem2-ineq-1}), $x\geq \ell e^{2(\ell+1)}$, and (\ref{lem2-ineq-2}), we have
 	\begin{align*}
 		\pi(x)-\pi(cx)
 		&> \left(\frac{1}{8(\ell+1)}-\frac{0.0008375}{2(\ell+1)}-\frac{0.0008375}{2(\ell+1)}\right)\frac{x}{\log x}\\
 		&\geq 0.1241625 \cdot \frac{1}{\ell+1}\frac{x}{\log x}.
 	\end{align*}
 	
 	\textbf{Case 2.} For $\max\{\ell e^{2(\ell+1)},7\cdot 10^{5}\}\leq x< 1.573\cdot 10^{9}$, there exist two integers $u,k$ with
 	$2\le u\le 6$ and  $10^3\le k<10^4$ such that $10^uk\le
 	x<10^u(k+1)$. With the help of a computer\footnote{In Case 2, we employed a Python program with the Python packages math and sympy to iterate over all integers $u$ and $k$ satisfying $10^{u}k\leq x<10^{u}(k+1)$, with the aim of finding the minimum value of the coefficient of $\frac{1}{\ell+1}$.}, we have
%%%%%%%%%%%%%%%%%%%%%%%%%%%
%%%%%%%%%%%%%%%%%%%%%%%%%%%
% Starting from here, there are many "with the help of (a) computer".    
% It is better to mention a more detail. For example, "with the help of Maple"
% It is even better to give some program (if it is not too complicated) here 
% or in Appendix.  From the second time of "with the help of ...", 
% we may not be necessary to give any more details.    
%%%%%%%%%%%%%%%%%%%%%%%%%%%
%%%%%%%%%%%%%%%%%%%%%%%%%%%
 	\begin{align*}
 		\big(\pi(x)-\pi(cx)\big)\frac{\log x}{x}
 		\geq &\Big( \pi(10^{u}k)-\pi \big(c \cdot 10^{u}(k+1) \big) \Big)\frac{\log \left( 10^{u}(k+1) \right)}{10^{u}(k+1)}\\
 		\geq &0.11604925\ldots \cdot \frac{1}{\ell+1},
 	\end{align*}
 	and the last equality holds if $\ell=8$, $u=6$, and $k=1001$.
 	
 	This completes the proof of Lemma \ref{lem-2}.
 \end{proof}

\end{document}
